\input{sol_pre.tex}
\input{q_2025-02-23_1.tex}
\begin{document}
\input{sol_head.tex}

\Savol{1}
$1111a-111b+11c+d+f=5000$. Faqat $a=5$ mos: $111b=555+11c+d+f\le555+99+17$ $\Rightarrow$ $b=6$, \ $11c+d+f=111$.\\
$c=9$, $d+f=12$; raqamlar turli ($5,6,9$ band): $\{d;f\}=\{4;8\}$. \ $ab-df+c=30-32+9=7$.
\Javob{C}

\Savol{2}
$4{,}(6)=\frac{14}3$, \ $\frac{14}3\cdot\frac{18}7=12$; \ $1{,}0(5)=\frac{19}{18}$: \ $9\cdot\frac{19}{18}-2{,}5=9{,}5-2{,}5=7$.
\Javob{B}

\Savol{3}
Unumdorliklar $p_1,p_2$ (töldiradi), $p_3,p_4$ (böşatadi):
\[p_1+p_2-p_3-p_4=\tfrac15,\quad p_1+p_2-p_3=\tfrac13\ \Rightarrow\ p_4=\tfrac2{15};\qquad p_1-p_3-p_4=\tfrac1{30}\ \Rightarrow\ p_1-p_3=\tfrac16.\]
Javob: $6$ soat.
\Javob{C}

\Savol{4}
$100\,000\cdot1{,}1\cdot0{,}85=93\,500$.
\Javob{A}

\Savol{5}
$x^2=2^{11\cdot24}$ $\Rightarrow$ $x=2^{132}$; \ $\sqrt[3]y=2^{45}$ $\Rightarrow$ $y=2^{135}$. \ $\dfrac xy=2^{-3}=\dfrac18$.
\Javob{B}

\Savol{6}
$\sqrt{15}+\sqrt{33}-\sqrt{22}-\sqrt{10}=\sqrt3\left(\sqrt5+\sqrt{11}\right)-\sqrt2\left(\sqrt{11}+\sqrt5\right)=\left(\sqrt{11}+\sqrt5\right)\left(\sqrt3-\sqrt2\right)$:
\[\frac{(11-5)\left(\sqrt3-\sqrt2\right)}{6\sqrt3-6\sqrt2}=1.\]
\Javob{A}

\Savol{7}
\[\frac{\sqrt2+\sqrt3+\sqrt5}{\left(\sqrt2+\sqrt3\right)^2-5}=\frac{\sqrt2+\sqrt3+\sqrt5}{2\sqrt6}=\frac{\left(\sqrt2+\sqrt3+\sqrt5\right)\sqrt6}{12}=\frac{3\sqrt2+2\sqrt3+\sqrt{30}}{12}.\]
\Javob{B}

\Savol{8}
$S_{28}=14(a_{14}+a_{15})=600$ $\Rightarrow$ $a_{15}=\frac{160}7$, \ $d=\frac{20}7$. \ $a_{28}=a_{14}+14d=20+40=60$.
\Javob{A}

\Savol{9}
$b_1\left(1+q+q^2\right)=14$, \ $b_1^2\left(1+q^2+q^4\right)=84$. \ $1+q^2+q^4=\left(1+q+q^2\right)\left(1-q+q^2\right)$ $\Rightarrow$ $b_1\left(1-q+q^2\right)=6$.\\
Ayirmasi: $2b_1q=8$, yiğindisi: $b_1\left(1+q^2\right)=10$ $\Rightarrow$ $\frac4q+4q=10$ $\Rightarrow$ $q=\frac12$ (kamayuvçi), $b_1=8$.
\[S_6=\frac{8\left(1-\frac1{64}\right)}{\frac12}=\frac{63}4=15\frac34.\]
\Javob{A}

\Savol{10}
$P$ nuqta $\frac\pi2$ burçakka mos. $\frac\pi2+3\approx4{,}57$: \ $\pi<4{,}57<\frac{3\pi}2\approx4{,}71$ — III çorak.
\Javob{B}

\Savol{11}
$a<c<0<b$: \ $|b-c|=b-c$, \ $|c-a|=c-a$, \ $|a-b|=b-a$; \ $\frac a{|a|}=\frac c{|c|}=-1$, $\frac b{|b|}=1$:
\[-(b-c)+(c-a)-(b-a)=2c-2b.\]
\Javob{D}

\Savol{12}
$4\cdot4^{x-1}\cdot3^{2x+1}=3\cdot36^x$, \ $12^{x+1}\cdot3^{x+2}=108\cdot36^x$; \ maxraj: $8\cdot36^x+2\cdot36^x=10\cdot36^x$:
\[\frac{-111\cdot36^x}{10\cdot36^x}=-11{,}1.\]
\Javob{C}

\Savol{13}
$t=\tg^22\alpha$: \ $\sin^22\alpha=\frac t{1+t}$, $\cos^22\alpha=\frac1{1+t}$:
\[\frac{t^3}{(1+t)(t-1)}-\frac1{t(1+t)(t-1)}=\frac{t^4-1}{t\left(t^2-1\right)}=t+\frac1t.\]
$\tg^22\alpha+\ctg^22\alpha+2=\left(\tg2\alpha+\ctg2\alpha\right)^2=\left(\dfrac2{\sin4\alpha}\right)^2=\dfrac4{\sin^24\alpha}$.
\Javob{C}

\Savol{14}
$3^x-3^{x-1}=2\cdot3^{x-1}$, \ $-2^{x+2}+2^{x-3}+2^{x-1}=2^{x-3}(-32+1+4)=-27\cdot2^{x-3}$:
\[2\cdot3^{x-1}\ge27\cdot\frac{2^{x-1}}4\ \Rightarrow\ \left(\frac32\right)^{x-1}\ge\frac{27}8\ \Rightarrow\ x\ge4.\]
\Javob{C}

\Savol{15}
$2023^{\log^2_{2023}2024}=\left(2023^{\log_{2023}2024}\right)^{\log_{2023}2024}=2024^{\log_{2023}2024}$ — birinçi ikki had qisqaradi.\\
$2023$ ta ildiz: $2^{1/2^{2023}}$, \ $\log_2\log_22^{2^{-2023}}=-2023$. \ Natija $0-(-2023)=2023$.
\Javob{D}

\Savol{16}
$x\ne-7$: \ $(x+2)^2=25$ $\Rightarrow$ $x=3$ yoki $x=-7$ (yaroqsiz). \ $1$ ta ildiz.
\Javob{A}

\Savol{17}
$(x-3)^2+3x=x^2-3x+9$. \ $t=\sqrt{x^2-3x+7}\ge0$: \ $t^2+2=t+22$ $\Rightarrow$ $t=5$.\\
$x^2-3x-18=0$, \ köpaytma $-18$.
\Javob{B}

\Savol{18}
$\left(\sqrt2x+1\right)^2\le0$ $\Rightarrow$ $x=-\frac1{\sqrt2}$ — bitta yeçim.
\Javob{A}

\Savol{19}
$\dfrac{10}{4-x}=-\dfrac{10}{x-4}$:
\[\frac{x^2-5x+64+10(x-6)}{(x-4)(x-6)}=\frac{(x+1)(x+4)}{(x-4)(x-6)}\le0\ \Rightarrow\ x\in[-4;\,-1]\cup(4;\,6).\]
Butunlar: $-4-3-2-1+5=-5$.
\Javob{D}

\Savol{20}
$E(f)=\left[-\sqrt{52};\,\sqrt{52}\right]$, \ $7<\sqrt{52}<8$: $1,\dots,7$ — $7$ ta.
\Javob{B}

\Savol{21}
Viyet: $2\cdot3=6$ — bajariladi; \ $a^2-2a-3=2+3$ $\Rightarrow$ $a^2-2a-8=0$ $\Rightarrow$ $a=4;\ -2$.
\Javob{C}

\Savol{22}
$g(x)=\dfrac{\ln\sin x}{\ln\cos x}$, \ $g\left(\frac\pi4\right)=1<\ln\pi$ $\Rightarrow$ $\frac\pi4$ atrofida $f=\ln\pi-g$, \ $f'=-g'$.
\[g'=\frac{\ctg x\ln\cos x+\tg x\ln\sin x}{\ln^2\cos x},\qquad x=\frac\pi4:\ L=\ln\frac{\sqrt2}2=-\frac12\ln2,\ \ g'=\frac{2L}{L^2}=-\frac4{\ln2}.\]
$f'\left(\frac\pi4\right)=\dfrac4{\ln2}=4\log_2e$.
\Javob{D}

\Savol{23}
$5-4x-x^2=-(x-1)(x+5)$: \ $\dfrac3{5-4x-x^2}=-\dfrac12\left(\dfrac1{x-1}-\dfrac1{x+5}\right)$.
\[-\frac12\ln\frac{x-1}{x+5}\Big|_2^4=-\frac12\left(\ln\frac13-\ln\frac17\right)=-\frac12\ln\frac73=\ln\sqrt{\frac37}.\]
\Javob{B}

\Savol{24}
Urinmalar tengligidan: $AE=p-BC=\dfrac{9+7+10}2-7=6$.
\Javob{B}

\Savol{25}
$AB=15\cdot0{,}6=9$, \ $AC=12$, \ $S=\frac12\cdot9\cdot12=54$.
\Javob{B}

\Savol{26}
$DE=EF$ va $\angle DEF=90^\circ$. $EFC$ uçburçakni $E$ atrofida $90^\circ$ ga burib, $F$ ni $D$ ga keltiramiz: $C$ nuqta $BD$ töğri çiziğiga tuşadi va $\angle BEC'=90^\circ$. Böyalgan soha — katetlari $20$ va $30$ bölgan töğri burçakli uçburçak:
\[S=\frac12\cdot20\cdot30=300.\]
\Javob{C}

\Savol{27}
$\triangle AMD$ da $\angle MAD=45^\circ$: \ $b=\angle FMA=45^\circ+\angle ADF$, \ $\tg\angle ADF=\frac26=\frac13$.\\
$\triangle CNE$ da $\angle NCE=45^\circ$: \ $a=180^\circ-45^\circ-\angle CEB$, \ $\tg\angle CEB=\frac63=2$.
\[a+b=180^\circ-\left(\angle CEB-\angle ADF\right),\qquad \tg\left(\angle CEB-\angle ADF\right)=\frac{2-\frac13}{1+\frac23}=1.\]
$a+b=180^\circ-45^\circ=135^\circ$.
\Javob{B}

\Savol{28}
$n=\frac{360^\circ}{45^\circ}=8$. \ $S=\frac12\cdot8\cdot R^2\sin45^\circ=4\cdot4\cdot\frac{\sqrt2}2=8\sqrt2$.
\Javob{D}

\Savol{29}
$BP$, $DQ$ — qirraga perpendikulyarlar: $AP=12\cos30^\circ=6\sqrt3$, $BP=6$; \ $CQ=4\sqrt3\cos60^\circ=2\sqrt3$, $DQ=6$.\\
$AQ=7\sqrt3-2\sqrt3=5\sqrt3$ $\Rightarrow$ $PQ=\sqrt3$. $BP$ va $DQ$ orasidagi burçak $120^\circ$:
\[BD^2=PQ^2+BP^2+DQ^2-2\cdot BP\cdot DQ\cos120^\circ=3+36+36+36=111.\]
\Javob{D}

\Savol{30}
$\angle BAO=30^\circ$ — içki çizilgan burçak, markaziy burçak $\angle BOx=60^\circ$. $B$ dan $Ox$ gaça: $6\sin60^\circ=3\sqrt3$.
\Javob{B}

\Savol{31}
$A\cap B=\varnothing$ $\Rightarrow$ $(A\cup B)'\cup A=B'$. \ Universal töplamda $10$ ta, $B$ da $3$ ta element: $n(B')=7$, \ $2^7=128$.
\Javob{A}

\Savol{32}
Toq raqamlar qoldiqlari ($\bmod3$): $0$: $\{3;9\}$, \ $1$: $\{1;7\}$, \ $2$: $\{5\}$. Qoldiqlar yiğindisi $3$ ga karrali bölişi kerak:
\[(0,0,0):\ 8;\quad(1,1,1):\ 8;\quad(2,2,2):\ 1;\quad(0,1,2)\ \text{tartibda}:\ 6\cdot2\cdot2\cdot1=24.\]
Jami $8+8+1+24=41$.
\Javob{C}

\newpage
\Savol{33}
Konus va silindr balandligi $h$: $2h=6$ $\Rightarrow$ $h=3$. Radius $r$, yasovçi $l=\sqrt{r^2+9}$.

\textbf{33.} $\pi rl=20\pi$ $\Rightarrow$ $r^2\left(r^2+9\right)=400$ $\Rightarrow$ $r^2=16$. \ $V=\pi r^2\cdot3+\frac13\pi r^2\cdot3=4\pi r^2=64\pi$ — \textbf{A}.

\textbf{34.} $\frac13\pi r^2\cdot3=81\pi$ $\Rightarrow$ $r=9$. \ $2\pi rh=54\pi$ — \textbf{C}.

\textbf{35.} $r^2+9=\frac{17}{16}r^2$ $\Rightarrow$ $r=12$. \ $2\pi\cdot12\cdot3=72\pi$ — \textbf{F}.
\Javob{33 — A, \ 34 — C, \ 35 — F}

\newpage
\Savol{36}
$x^4-4x^3=\left(x^2-2x\right)^2-4x^2$, şuning uçun
\[\left(x^2-2x\right)^2-8\left(x^2-2x+1\right)=0\ \Rightarrow\ \left(x^2-2x\right)^2=8(x-1)^2.\]
$w=(x-1)^2$, \ $x^2-2x=w-1$: \ $(w-1)^2=8w$ $\Rightarrow$ $w^2-10w+1=0$ $\Rightarrow$ $w=5\pm2\sqrt6=\left(\sqrt3\pm\sqrt2\right)^2$.
\[x=1\pm\left(\sqrt3+\sqrt2\right),\qquad x=1\pm\left(\sqrt3-\sqrt2\right).\]
a) $4$ ta haqiqiy ildiz.

b) Faqat $1-\sqrt3-\sqrt2<0$, qolganlari musbat:
\[\left(1+\sqrt3+\sqrt2\right)+\left(\sqrt3+\sqrt2-1\right)+\left(1+\sqrt3-\sqrt2\right)+\left(1-\sqrt3+\sqrt2\right)=2+2\sqrt2+2\sqrt3.\]
\Javob{a) 4; \quad b) $2+2\sqrt2+2\sqrt3$}

\newpage
\Savol{37}
$p=\tg x$, $q=\tg\frac y2$: \ $p+q=\dfrac4{\sqrt3}$, \ $\dfrac{p+q}{pq}=\dfrac4{\sqrt3}$ $\Rightarrow$ $pq=1$.
\[t^2-\frac4{\sqrt3}t+1=0\ \Rightarrow\ t=\sqrt3\ \text{yoki}\ t=\frac1{\sqrt3}.\]
a) $\tg y=\dfrac{2q}{1-q^2}$: \ $q=\sqrt3$: $-\sqrt3$; \ $q=\frac1{\sqrt3}$: $\sqrt3$. Eng kiçigi $-\sqrt3$.

b) $\tg x=\frac1{\sqrt3}$ yoki $\sqrt3$: eng kiçik musbat $x=\dfrac\pi6$.
\Javob{a) $-\sqrt3$; \quad b) $\frac\pi6$}

\newpage
\Savol{38}
$x$ ni $\frac1x$ ga almaştiramiz: $f\left(\frac1x\right)+\dfrac{1-x}{2+x}f(x)=\dfrac{3x}{1-x}$. Bundan $f\left(\frac1x\right)$ ni birinçi tenglikka qöyamiz:
\[f(x)\left(1+\frac{(x-1)^2}{(2x+1)(x+2)}\right)=\frac3{x-1}+\frac{3x}{2x+1}\ \Rightarrow\ f(x)\cdot\frac{3\left(x^2+x+1\right)}{(2x+1)(x+2)}=\frac{3\left(x^2+x+1\right)}{(x-1)(2x+1)}.\]
$f(x)=\dfrac{x+2}{x-1}$.

a) $f(2)=4$. \quad b) $f(x)=0$ $\Rightarrow$ $x=-2$.
\Javob{a) 4; \quad b) $-2$}

\newpage
\Savol{39}
Grafikda «tepalar» — lokal maksimumlar, «çuqurliklar» — lokal minimumlar. Silliq grafikda bu nuqtalarda $f'(x)=0$.

a) Lokal maksimumlar: $3$ ta.

b) $3$ ta maksimum va $3$ ta minimum: $f'(x)=0$ ning $6$ ta ildizi.
\Javob{a) 3; \quad b) 6}

\newpage
\Savol{40}
$y_1(2)=\frac{64}{16}=4$, \ $y_1'=\frac{3(x+2)^2}{16}$, \ $y_1'(2)=3$: \ $y_2=3(x-2)+4=3x-2$.

a) $k=3$, $b=-2$: \ $kb^2=12$.

b) Soha $Oy$, $Ox$, $y_1$ va $y_2$ bilan çegaralangan; $y_2$ $Ox$ ni $x=\frac23$ da kesadi:
\[S=\int_0^2\frac{(x+2)^3}{16}dx-\frac12\cdot\left(2-\frac23\right)\cdot4=\frac{(x+2)^4}{64}\Big|_0^2-\frac83=\frac{240}{64}-\frac83=\frac{13}{12}.\]
\Javob{a) 12; \quad b) $\frac{13}{12}$}

\newpage
\Savol{41}
$\triangle ABH\sim\triangle BCH\sim\triangle ACB$ — içki aylana radiuslari gipotenuzalarga proporsional:
\[\frac{AB}{BC}=\frac34,\qquad r^2=r_1^2+r_2^2\ \left(AB^2+BC^2=AC^2\right)\ \Rightarrow\ r=5.\]
$AB=3k$, $BC=4k$, $AC=5k$: \ $r=\frac{3k+4k-5k}2=k=5$ $\Rightarrow$ $AB=15$, $BC=20$, $AC=25$.

a) $BH=\dfrac{AB\cdot BC}{AC}=\dfrac{15\cdot20}{25}=12$.

b) $S=\frac12\cdot15\cdot20=150$.
\Javob{a) 12; \quad b) 150}

\newpage
\Savol{42}
Içki aylana mavjud $\Rightarrow$ qarama-qarşi tomonlar yiğindisi teng, demak teng tomonlar qöşni: $AB=BC=6$, $CD=DA=8$ (deltoid). $BD$ — simmetriya öqi, $\angle A=\angle C$; taşqi aylana: $\angle A+\angle C=180^\circ$ $\Rightarrow$ $\angle A=\angle C=90^\circ$.
\fig{\begin{tikzpicture}[scale=0.45,font=\small]
\coordinate (B) at (-5,0); \coordinate (D) at (5,0); \coordinate (A) at (-1.4,4.8); \coordinate (C) at (-1.4,-4.8); \coordinate (O) at (0,0); \coordinate (I) at (-0.714,0); \coordinate (K) at (-1.4,0); \coordinate (T) at (-2.971,2.743);
\draw (O) circle (5); \draw (I) circle (3.429); \draw (A)--(B)--(C)--(D)--cycle;
\draw[axs] (B)--(D); \draw[ax] (A)--(C); \draw[ax] (I)--(T);
\fill (O) circle (3pt); \fill (I) circle (3pt);
\node[above] at (A) {$A$}; \node[left] at (B) {$B$}; \node[below] at (C) {$C$}; \node[right] at (D) {$D$};
\node[below right] at (O) {$O$}; \node[below] at (I) {$I$}; \node[aux,above] at (T) {$T$};
\end{tikzpicture}}
$BD=\sqrt{36+64}=10$ (diametr, $R=5$).

a) $AC=2\cdot\dfrac{AB\cdot AD}{BD}=2\cdot\dfrac{48}{10}=9{,}6$ — kiçik diagonal.

b) $S=2\cdot\frac12\cdot6\cdot8=48$, \ $r=\frac Sp=\frac{48}{14}=\frac{24}7$. Içki markaz $I$ — $BD$ da; $IT\perp AB$, $\sin\angle ABD=\frac{AD}{BD}=\frac45$:
\[BI=\frac r{\sin\angle ABD}=\frac{24/7}{4/5}=\frac{30}7,\qquad OI=BO-BI=5-\frac{30}7=\frac57.\]
\Javob{a) $9{,}6$; \quad b) $\frac57$}

\newpage
\Savol{43}
\fig{\begin{tikzpicture}[scale=1.1,font=\small]
\coordinate (A) at (-1.823,1.823); \coordinate (B) at (-1,0); \coordinate (C) at (0,-1.732); \coordinate (D) at (1,0); \coordinate (E) at (1.823,1.823); \coordinate (F) at (0,2.646); \coordinate (P) at (0,0);
\fill[okgreen!30] (B)--(F)--(D)--cycle;
\draw (A)--(B)--(C)--(D)--(E)--(F)--cycle; \draw[axs] (B)--(F)--(D)--cycle; \draw[ax] (F)--(P);
\node[left] at (A) {$A$}; \node[left] at (B) {$B$}; \node[below] at (C) {$C$}; \node[right] at (D) {$D$}; \node[right] at (E) {$E$}; \node[above] at (F) {$F$}; \node[aux,below] at (P) {$P$};
\end{tikzpicture}}
$\triangle BCD$: $BC=CD=2$, $\angle C=60^\circ$ $\Rightarrow$ $BD=2$. \ $\triangle FAB$, $\triangle DEF$ — teng yonli töğri burçakli: $BF=DF=2\sqrt2$.

a) $BFD$ teng yonli, $FP$ — balandlik: $FP=\sqrt{8-1}=\sqrt7$, \ $S_{BFD}=\frac12\cdot2\cdot\sqrt7=\sqrt7$.

b) Tomon $2$ bölganda: $S=S_{FAB}+S_{DEF}+S_{BCD}+S_{BFD}=2+2+\sqrt3+\sqrt7$.\\
Tomon $a$ da yuz $\left(\frac a2\right)^2$ marta özgaradi: $a^2=\dfrac{100}{4+\sqrt7+\sqrt3}$:
\[S=\frac{a^2}4\left(4+\sqrt3+\sqrt7\right)=25.\]
\Javob{a) $\sqrt7$; \quad b) 25}

\newpage
\Savol{44}
Öq kesimi: teng yonli uçburçak, asosdagi burçaklar $75^\circ$, uçidagi burçak $30^\circ$; şar katta doirasi unga taşqi çizilgan.
\fig{\begin{tikzpicture}[scale=0.32,font=\small]
\coordinate (S) at (0,6); \coordinate (A) at (-3,-5.196); \coordinate (B) at (3,-5.196); \coordinate (O) at (0,0); \coordinate (H) at (0,-5.196);
\draw (O) circle (6); \draw (A)--(S)--(B)--cycle; \draw[ax] (S)--(H); \draw[axs] (O)--(B);
\fill (O) circle (4pt);
\node[above] at (S) {$S$}; \node[below left] at (A) {$A$}; \node[below right] at (B) {$B$}; \node[left] at (O) {$O$}; \node[aux,below] at (H) {$H$};
\node at (-1.6,-4.4) {\scriptsize$75^\circ$};
\end{tikzpicture}}
a) $l=\dfrac r{\cos75^\circ}$, \ $\dfrac1{\cos75^\circ}=\dfrac4{\sqrt6-\sqrt2}=\sqrt6+\sqrt2$, \ $\sqrt{2+\sqrt3}=\dfrac{\sqrt6+\sqrt2}2$:
\[\pi rl=\pi r^2\left(\sqrt6+\sqrt2\right)=18\cdot\frac{\sqrt6+\sqrt2}2\pi\ \Rightarrow\ r^2=9,\ r=3.\]
b) Sinuslar teoremasi: $2R=\dfrac{AB}{\sin30^\circ}=\dfrac6{1/2}=12$, \ $R=6$.
\[V=\frac43\pi R^3=\frac43\cdot3\cdot216=864.\]
\Javob{a) 3; \quad b) 864}

\newpage
\Savol{45}
$r=20$ cm $=2$ dm, \ uzunlik $L=200$ cm $=20$ dm.

a) Yarim silindr yon sirti va ikkita yarim doira:
\[\pi rL+2\cdot\frac12\pi r^2=3\cdot2\cdot20+3\cdot4=120+12=132\ \text{dm}^2.\]
b) $V=\frac12\pi r^2L=\frac12\cdot3\cdot4\cdot20=120$ dm$^3$ $=120$ litr.
\Javob{a) 132; \quad b) 120}

\ShowJavoblar
\end{document}
